\section{Introduction}
\label{sec:introduction}

\subsection{Background: forecasting things that never resolve}

Prediction markets are one of the most successful practical tools for aggregating
dispersed information. By letting people bet on the outcome of an event and paying
them once the truth is revealed, markets such as the Iowa Electronic Markets and,
more recently, Polymarket~\parencite{polymarket} produce forecasts that are
routinely competitive with polls and expert panels~\parencite{berg2008results}.
The engine underneath is a \emph{proper scoring rule}: a payment scheme that makes
honestly reporting one's true belief the reward-maximising action, provided the
outcome is eventually observed~\parencite{brier1950verification,good1952rational,gneiting2007strictly}.

The catch is the phrase ``provided the outcome is eventually observed''. Many of the
questions decision-makers care about most are \emph{unverifiable}: the ground truth
is hard, expensive, or simply impossible to access. What was the causal effect of a
policy that was actually enacted, relative to the counterfactual that was not? Is a
given piece of online content genuinely misinformation? What will the long-run
social consequence of a choice be? In each case people must still pool scattered
information and commit to a ``best guess'', yet there is no future observation to
pay against. The standard prediction-market lever is missing.

\subsection{Context: previous approaches}

Two research traditions speak to this problem but neither solves it cleanly. The
first is \emph{Aumann's agreement theorem} and its descendants~\parencite{aumann1976agreeing,kong2022aggregation}:
rational Bayesians who share a common prior and repeatedly exchange beliefs must
eventually agree, and --- when their signals are informational
substitutes~\parencite{chen2016informational} --- that consensus aggregates all of
their private information. This machinery aggregates beliefs but assumes agents
\emph{report honestly}; it provides no \emph{incentive} to reveal a privately held
signal. The second tradition is \emph{information elicitation without verification},
which includes peer prediction~\parencite{miller2005eliciting}, the Bayesian truth
serum~\parencite{prelec2004bayesian}, and output-agreement
mechanisms~\parencite{waggoner2014output}. These provide incentives without an
outcome, but they typically target \emph{elicitation} from a few agents rather than
efficient \emph{aggregation} from a large pool, and many of them admit
uninformative equilibria in which everyone reports the same uninformative thing.

\subsection{The problem: the gap}

There is therefore a gap between these literatures. \emph{However}, the agreement
literature aggregates information but does not incentivise truthful revelation.
\emph{Despite} the elicitation literature providing incentives without ground
truth, it does not aggregate efficiently and is fragile to uninformative equilibria.
What is missing is a single mechanism that simultaneously (i) elicits truthful
private beliefs, (ii) aggregates them into one consensus forecast, and (iii)
requires \emph{no} access to the outcome --- all under roughly the same assumptions
that make ordinary prediction markets work.

\subsection{The solution: a well-informed reference replaces the outcome}

\textcite{srinivasan2025selfresolving} close this gap with a
\emph{self-resolving prediction market}. Agents report sequentially; after each
report the market terminates with probability $\alpha$; and instead of paying agents
against the outcome, the mechanism pays them the cross-entropy of their report
against the \emph{final (reference) agent}, who --- having observed the full market
history --- is the best-informed participant. The key insight is one of
\emph{informational smallness}: if the reference agent has access to enough
conditionally independent signals (``informational substitutes''), then no single
earlier agent can meaningfully move the reference's posterior, so the reference
behaves like a noisy but unbiased proxy for the ground truth. Paying against this
proxy recovers the truthful incentives of a standard market \emph{without ever
observing the outcome}.

\subsection{The proof: summary of the evaluation}

The original paper establishes that truthful reporting is a strict perfect Bayesian
equilibrium under a standard set of assumptions, but it offers no empirical
evaluation: the mechanism is never simulated and never run on data. This paper
supplies that missing evaluation. We build an agent-based simulator in which the
paper's assumptions hold by construction, and we ask two questions. \textbf{RQ1
(incentives, simulation):} under strategic deviation, when is truthful reporting
actually a best response, and how does this depend on how well-informed the
reference is? We find that a focal agent's profit-maximising deviation
$\gamma^\star$ converges to the truthful value $1$ as the reference's number of
informational substitutes $k$ grows --- pegged at gross over-reporting for $k \le 8$,
$\gamma^\star \approx 1.2$ at $k = 32$, and exactly $1.0$ for $k \ge 64$ ---
operationalising the paper's Theorems~1--3. \textbf{RQ2 (robustness, empirical):}
on roughly $90$ resolved Polymarket markets we pay each price path both
self-resolvingly (against a late reference price) and verifiably (against the
realised outcome), and find that the two schemes agree almost perfectly at a
genuinely late reference (Pearson $\approx 1.00$, payout ratio $\approx 0.95$) and
degrade gracefully as the reference is moved earlier --- exactly tracing out the
outcome-leakage channel the mechanism depends on.

\subsection{Contributions}

Concretely, this project contributes:
\begin{enumerate}
    \item \textbf{An open-source agent-based simulator} (\texttt{srpm}) of the
    self-resolving prediction market, with configurable agents, signal structures,
    scoring rules, and strategies, in which the paper's rationality, common-prior,
    distinct-signal, and conditional-independence assumptions hold by construction
    (Section~\ref{sec:methods}).
    \item \textbf{An incentive (best-response) study (RQ1)} that estimates a focal
    agent's profit-maximising deviation by Monte Carlo and shows it converges to
    truthful reporting as the reference's informational substitutes grow, together
    with an exact reproduction of the paper's analytic minimum-substitutes bound
    $k_{\min}$ (Sections~\ref{sec:methods}--\ref{sec:experiments}).
    \item \textbf{Supporting synthetic studies} confirming that under truthful play
    the market reproduces the full-information Bayesian posterior to machine
    precision, that a few near-oracle references pin the aggregate, and that the
    mechanism pays almost nothing in uninformative equilibria
    (Section~\ref{sec:experiments}).
    \item \textbf{A counterfactual empirical study (RQ2)} on real resolved
    Polymarket data that quantifies how closely self-resolving payouts track the
    verifiable payouts they are meant to replace, and isolates outcome leakage as
    the mechanism behind the residual gap (Section~\ref{sec:experiments}).
\end{enumerate}

The remainder of the paper is organised as follows. Section~\ref{sec:related}
reviews related work; Section~\ref{sec:preliminaries} sets up the model, the
mechanism, and its assumptions; Section~\ref{sec:methods} describes our simulator
and experimental design; Section~\ref{sec:experiments} presents the results; and
Section~\ref{sec:conclusion} discusses limitations and concludes.
